% WP3 PRL figure callout and caption.
% Scope: minimal BTK-like observable slice, not full self-consistent BTK.

Figure~3 converts the minimal pairing-screening result into a directly observable conductance diagnostic. The three representative traces are chosen to separate response families rather than to maximize a single peak height. The ordinary same-sign $s$-wave channel gives a broad zero-bias enhancement without a resolved finite-bias splitting. The $s_{\pm}$ candidate remains weak in the transparent limit but develops a strong zero-bias response only when intervalley mixing is turned on, consistent with a mixing-assisted frustration channel. By contrast, the chiral channel already carries a finite-bias-dominated subgap response at zero intervalley mixing, with the representative trace peaking near $eV=-0.026$ and a positive-negative peak separation of about $0.052$ in the present energy units. This separation is not used as a unique pairing identifier. It supplies the observable layer required to distinguish a phase-sensitive finite-bias response from the smooth null-model background.

\begin{figure}[t]
  \centering
  % \includegraphics[width=0.95\linewidth]{figures/wp3_conductance_fingerprints.pdf}
  \caption{\label{fig:wp3-fingerprints}
  Minimal conductance fingerprints from the pairing-screening kernel.
  Representative BTK-like conductance traces for ordinary same-sign $s$-wave, intervalley-sensitive $s_{\pm}$, and chiral pairing channels are shown on the same voltage axis.
  The $s$-wave trace shows a broad zero-bias enhancement and no resolved finite-bias splitting.
  The $s_{\pm}$ response becomes strong only when intervalley mixing is present, consistent with mixing-assisted sign frustration.
  The chiral trace is finite-bias dominated already at zero intervalley mixing; in the selected slice the strongest peak is near $eV=-0.026$, with a positive-negative peak separation of approximately $0.052$.
  These curves are observable-level fingerprints of the minimal model and should not be read as a full microscopic BTK proof or a unique order-parameter identification.}
\end{figure}
